% !PRINTED 558
\weberpage{558}
%<<CONT
tanten 1 hat. Um entgegengesetzte Classen zu definiren, nehmen wir einen Repräsentanten $\varphi$ einer Classe $A$ und eine durch $\varphi$ theilbare Zahl $\alpha$. Ist dann $\alpha = \varphi\chi$, so ist $\chi$ ein Repräsentant der Classe $A^{-1}$, und es ist $AA^{-1} = O$.

Ist $A$ eine beliebige Classe, und $h$ die Classenzahl, so ist immer

\[ A^h = O, \]

und wenn $k$ die kleinste positive Zahl ist, die der Bedingung

\[ A^k = O \]

genügt, so ist $k$ ein Theiler von $h$. Daraus ergiebt sich der Satz:

3.~\gesperrt{Jedes Functional $\varphi$ in $\Omega$ gehört zu einem bestimmten Exponenten $k$, der ein Theiler der Classenzahl $h$ ist, so dass $\varphi^k$ associirt ist mit einer Zahl in $\Omega$.}

\begin{center}
\rule{4cm}{0.4pt}
\end{center}
